\section{Możliwości dalszego rozwoju}
\label{sec:rozwoj}

Wdrożony system stanowi działający \textit{proof of concept} architektury do przetwarzania i~archiwizacji danych
rynkowych w~czasie rzeczywistym w~środowisku \textit{on-premises}.
Z~uwagi na mikrousługowy charakter rozwiązania oraz silnie operacyjny kontekst przetwarzania strumieniowego,
istnieje szereg naturalnych kierunków rozwoju, które pozwoliłyby zwiększyć niezawodność, ułatwić diagnozę
problemów oraz rozszerzyć funkcjonalność analityczną.
Poniżej przedstawiono wybrane propozycje.

\subsection{Centralizacja logów}

W~obecnej postaci każdy mikroserwis zapisuje logi niezależnie, co utrudnia szybkie odtworzenie przebiegu zdarzeń
w~sytuacji awarii lub rozbieżności danych.
Wprowadzenie centralizacji logów (np.\ poprzez standaryzację logów w~formacie JSON, dodanie identyfikatora korelacji
oraz wysyłkę do wspólnego backendu typu Loki/ELK) umożliwiłoby analizę całego systemu w~jednym miejscu.
Taka zmiana skraca czas diagnostyki oraz ułatwia analizę zależności między komponentami.

\subsection[Metryki i dashboardy (Prometheus/Grafana)]{Metryki i~dashboardy (Prometheus/Grafana)}

Kolejnym krokiem jest ekspozycja metryk technicznych i~domenowych, takich jak: opóźnienie od odbioru zdarzenia do
zapisania w~bazie, procent błędów walidacji, opóźnienie konsumentów JetStream, czasy odpowiedzi API GraphQL czy wydajność
procesów archiwizacji.
Wysyłanie tych metryk do systemu typu Prometheus oraz budowa dashboardów w~Grafanie pozwoliłaby nie tylko
monitorować bieżący stan, ale również planować zasoby na podstawie trendów i~wykrywać degradację wydajności zanim
stanie się odczuwalna dla użytkownika~\cite{prometheus_grafana}.

\subsection{Śledzenie rozproszone (distributed tracing)}

W~systemie sterowanym zdarzeniami pojedyncze zdarzenie przechodzi przez wiele komponentów, co utrudnia śledzenie jego
przebiegu.
Wdrożenie śledzenia rozproszonego (np.\ OpenTelemetry) wraz z~propagacją kontekstu (\texttt{trace\_id}) między
mikroserwisami pozwoliłoby na wizualizację pełnej ścieżki przetwarzania i~precyzyjne wskazanie wąskich gardeł~\cite{open_telemetry}.
Jest to szczególnie istotne przy skalowaniu systemu oraz przy analizie incydentów trudnych do odtworzenia.

\subsection[Alerty operacyjne i powiadomienia o stanie systemu]{Alerty operacyjne i~powiadomienia o~stanie systemu}

W~projekcie istnieje bot Telegram wykorzystywany do powiadomień o~zdarzeniach rynkowych.
Naturalnym rozszerzeniem jest dodanie warstwy alertów operacyjnych, które informowałyby o~problemach działania
systemu: braku nowych danych, wysokim poziomie błędów, rosnącym opóźnieniu konsumentów, przekroczeniu progów
zużycia zasobów czy nieudanych migracjach danych do MinIO\@.
Alerty oparte o~metryki i~logi pozwoliłyby reagować szybciej i~skrócić czas przestojów.

\subsection[Rozszerzenie API: dostęp do danych z MinIO oraz federacja GraphQL]{Rozszerzenie API: dostęp do danych z~MinIO oraz federacja GraphQL}

Obecne API GraphQL udostępnia dane z~warstwy gorącej (MongoDB).
Dodanie zapytań do danych historycznych i~analitycznych przechowywanych w~MinIO umożliwiłoby obsługę szerszego
zakresu przypadków użycia, takich jak pobieranie agregatów, dłuższych okien czasowych oraz danych do treningu
modeli.

W~praktyce wygodnym podejściem jest zastosowanie federacji GraphQL, czyli mechanizmu pozwalającego łączyć wiele
niezależnych schematów i~serwerów GraphQL w~jeden spójny endpoint.
Federacja działa jak \textit{GraphQL dla endpointów GraphQL}: klient widzi jedno API pod jednym adresem, natomiast bramka
\textit{gateway} rozdziela zapytania do właściwych podserwisów i~składa odpowiedź w~całość~\cite{apollo_federation_architecture}.
Takie rozwiązanie upraszcza rozwój i~wdrażanie, ponieważ poszczególne części API mogą ewoluować niezależnie.

\subsection{Bardziej zaawansowany moduł analityczny}

Aktualna warstwa analityczna ma charakter demonstracyjny i~koncentruje się na podstawowej agregacji.
Kolejnym krokiem mogłoby być rozszerzenie zakresu obliczeń o~cechy realnie wykorzystywane w~analizie
ilościowej, takie jak zmienność w~oknach kroczących, miary trendu i~momentum, wykrywanie anomalii czy proste
sygnały zdarzeniowe.

\subsection[Modele predykcyjne i pipeline uczenia maszynowego]{Modele predykcyjne i~pipeline uczenia maszynowego}

System dostarcza obecnie dane oczyszczone i~zagregowane, co stanowi dobry fundament do budowy modeli predykcyjnych.
Dalszy rozwój mógłby obejmować przygotowanie pełnego potoku ML: definicję etykiet, walidację czasową,
rejestrowanie eksperymentów oraz wersjonowanie modeli.
W~wariancie produkcyjnym można rozważyć osobny serwis publikujący predykcje jako kolejne zdarzenia w
NATS, co pozwoliłoby wykorzystywać je w~powiadomieniach lub dalszych procesach ETL\@.

\subsection[Dalsza optymalizacja i skalowanie do wszystkich symboli]{Dalsza optymalizacja i~skalowanie do wszystkich symboli}

Aktualnie system obsługuje ograniczoną liczbę symboli, co upraszcza przetwarzanie i~zmniejsza wymagania zasobowe.
Skalowanie do pełnej listy symboli dostępnych na giełdzie wymagałoby poprawy mechanizmów \textit{backpressure},
fragmentacji strumieni i~konsumentów (np.\ podział po symbolu lub klasie instrumentów), grupowania zapisów do baz oraz
lepszego ograniczania kosztów obliczeniowych w~ETL\@.
Taki kierunek rozwoju jest istotny, jeśli celem systemu miałoby być pełne odwzorowanie rynku w~czasie rzeczywistym.
